\documentclass[10pt,a4paper]{amsart}
\usepackage[margin=19mm]{geometry}
\usepackage{amsmath,amssymb,mathtools}
\usepackage[scr=rsfs,cal=euler]{mathalfa}
\usepackage{xfrac}
\usepackage[unicode,hidelinks,bookmarks=false]{hyperref}
\newcommand{\ie}{i.e.\ }
\newcommand{\cf}{cf.\ }
\newcommand{\resp}{respectively\ }
\newcommand{\Subsection}{\subsection}
\providecommand{\nobreakdash}{\nobreak\hbox{-}}
\setlength{\emergencystretch}{2em}
\multlinegap0pt
\usepackage{enumerate}
\usepackage{amssymb}
\usepackage{tikz}
\usepackage{xcolor}
\newtheorem{theorem}{Theorem}
\newtheorem{proposition}{Proposition}
\def\d{\mathrm{d}}
\title{A lower bound for the first non-zero basic eigenvalue on a singular Riemannian foliation}
\author{Bach Tran}
\date{Source: arXiv:2602.17501v3 (CC BY 4.0)}
\begin{document}
\maketitle

\noindent\emph{Source and licence.} A lower bound for the first non-zero basic eigenvalue on a singular Riemannian foliation. Version arXiv:2602.17501v3 (CC BY 4.0), distributed by arXiv under the Creative Commons Attribution 4.0 International licence, http://creativecommons.org/licenses/by/4.0/, as recorded in the arXiv licence metadata for this exact version and shown on the abstract page https://arxiv.org/abs/2602.17501v3. Full licence notice: \texttt{licenses/arxiv-2602.17501-CC-BY-4.0.txt}.\par\smallskip

\noindent\textit{Editorial scope.} Counted unit: the article's theorem main1 (source0.tex lines 168--173). The article's complete original proof of it is delivered with it (source0.tex lines 476--530, one \texttt{proof} environment of 55 lines, which the article places away from the statement). No mathematics is written, repaired or re-derived: the statement and the proof are the article's own lines, in the article's own notation, and no retained line is substituted or re-flowed.  Retained uncounted as the source-local context the proof needs: lines 430--432; lines 449--459.  Nothing was omitted from any retained span, so the package carries no omission record.  No retained line contains a citation, so the wrapper carries no bibliography and no bibliography entry.  No retained line contains a figure environment, an image-inclusion command or drawing code, so no image is delivered and the package declares no asset dependency.  Editorial formatting only, in addition: an AMS \texttt{amsart} preamble that restates the article's own declarations and macros used by the retained lines, namely the environments \texttt{theorem}, \texttt{proposition} on separate counters, together with the standard packages those lines need.  The retained environments print in this document as Theorem 1, Proposition 1 in order of appearance, whereas the article prints the same items with the numbering of its own full document.  The complete native source of the article as distributed in the arXiv e-print for this exact version is bundled unchanged as \texttt{sources/194966/source0.tex}.
    \begin{align}\label{uia}
        v_{\varepsilon}=\frac{u-\frac{1-k}{2}}{(1+\varepsilon)\frac{1+k}{2}}.
    \end{align}

\begin{proposition}\label{oo}
With the setting above, we have that \begin{align}\label{zzzzz}
    \vert\nabla\theta_\varepsilon\vert^2\leq \lambda(1+a_\varepsilon\psi(\theta_\varepsilon)),
\end{align}where $\psi:\left[-\frac{\pi}{2},\frac{\pi}{2}\right]\to\mathbb{R}$ is the function given by
\begin{align*}
        \left\{\begin{array}{lc}
             \psi(\theta)=\frac{4}{\pi}(\theta\sec^2\theta+\tan\theta)-2\tan\theta\cdot\sec\theta,& \theta\in\left(-\frac{\pi}{2},\frac{\pi}{2}\right) \\
             \psi\left(\frac{\pi}{2}\right)=1,\text{ }\psi\left(-\frac{\pi}{2}\right)=-1& 
        \end{array}\right.
    \end{align*}
\end{proposition}

        \begin{theorem}\label{main1}

    Let $M^n$ be a compact Riemannian manifold with non-negative Ricci curvature and $\mathcal{F}$ a singular Riemannian foliation of $M$ with closed leaves and basic mean curvature. Then we have a lower bound for the first non-zero basic eigenvalue $\lambda_1^B$ of the Laplacian, given by \begin{align}
        \lambda_1^B\geq \frac{\pi^2}{d_{M/\mathcal{F}}^2}.
    \end{align}
\end{theorem}

        \begin{proof}[Proof of Theorem \ref{main1}]

        Let $u\in C^\infty_B(M)$ be a normalized basic eigenfunction corresponding to the eigenvalue $\lambda_1^B$, namely, a smooth basic function $u$ on $M$ satisfying $\Delta u=-\lambda_1^B u$ and $1=\max_{x\in M} u>\min _{x\in M}u\geq -1$. 

        For any $\varepsilon>0$, let $\theta_\varepsilon=\sin^{-1}v_\varepsilon$ where $v_\varepsilon$ is defined as in \ref{uia}. Notice that $\theta_\varepsilon$ is constant on each leaf of $\mathcal{F}$, namely, $\theta_\varepsilon$ is basic.
        
        Apply Proposition \ref{oo} to $\theta_\varepsilon$ and we get that
        \begin{align}\label{uiia}
            \vert\nabla\theta_\varepsilon\vert^2\leq\lambda_1^B(1+a_\varepsilon\psi(\theta_\varepsilon)),
        \end{align} and hence
        \begin{align}
            \label{aaaaa}\frac{\vert\nabla\theta_\varepsilon\vert}{\sqrt{1+a_\varepsilon\psi(\theta_\varepsilon)}}\leq\sqrt{\lambda_1^B}.
        \end{align}

        Let $x\in M$ such that $\theta_\varepsilon(x)=-\frac{\pi}{2}+\delta$.

        
        Choose $y$ from the level set $A=\theta_\varepsilon^{-1}(\frac{\pi}{2}-\delta)$ (which is compact and is a union of leaves of $\mathcal{F}$) minimizing the distance between $x$ and $A$ and let $\gamma
        $ be the geodesic joining $x$ and $y$. That means
        \begin{align}
            L(\gamma)=\d(x,y)=\d(x,A).
        \end{align}

        This geodesic $\gamma$ is also the shortest geodesic from $x$ to $L_y\subset A$, which means $d(x,A)=d(x,L_y)$.

        Since leaves of $\mathcal{F}$ are closed, the global equidistance of $\mathcal{F}$ indicates that $\d(x,L_y)=d(L_x,L_y)$ and hence
        \begin{align}
            L(\gamma)=\d(L_x,L_y)\leq\sup_{p,q\in M}\d(L_p,L_q)=d_{M/\mathcal{F}}
        \end{align}

        In \ref{aaaaa}, take the integral along $\gamma$, we get that
        \begin{align}\label{sf}
            \d_{M/\mathcal{F}}\sqrt{\lambda_1^B}\geq L(\gamma)\sqrt{\lambda_1^B}&=\int_\gamma{\sqrt{\lambda_1^B }\d t}\geq\int_\gamma\frac{\vert\nabla\theta_\varepsilon\vert}{\sqrt{1+a_\varepsilon\psi(\theta_\varepsilon)}} \d t\geq\left\vert\int_\gamma{\frac{\nabla\theta_\varepsilon\cdot \d r}{\sqrt{1+a_\varepsilon\psi(\theta_\varepsilon)}}}\right\vert.
        \end{align}

        Because $\psi$ is an odd function, we see that
        \begin{align}\label{sg}
            \left\vert\int_\gamma{\frac{\nabla\theta_\varepsilon\cdot \d r}{\sqrt{1+a_\varepsilon\psi(\theta_\varepsilon)}}}\right\vert&=\int_{-\frac{\pi}{2}+\delta}^{\frac{\pi}{2}-\delta}{\frac{\d\theta}{\sqrt{1+a_\varepsilon\psi(\theta)}}}\nonumber\\&=\int_{0}^{\frac{\pi}{2}-\delta
            }\left(\frac{1}{\sqrt{1+a_\varepsilon\psi(\theta)}}+\frac{1}{\sqrt{1-a_\varepsilon\psi(\theta)}}\right)\d\theta\nonumber\\&=2\int_0^{\frac{\pi}{2}-\delta}\left(1+\sum_{k=1}^{\infty}{\frac{(4k-1)!!}{(4k)!!}a_\varepsilon^{2k}\psi^{2k}(\theta)}\right)\d\theta\\&\geq 2\left(\frac{\pi}{2}-\delta\right)+2\frac{3!!}{4!!}a_\varepsilon^2\int_0^{\frac{\pi}{2}-\delta}\psi^2(\theta)\d\theta\nonumber\\&=\pi-2\delta+\frac{3}{4}a_\varepsilon^2\int_0^{\frac{\pi}{2}-\delta}\psi^2(\theta)d\theta\nonumber.
        \end{align}where ``$!!$'' is for double factorial.
        
        Finally, by letting $\varepsilon\to 0$, which also makes $\delta\to 0$, \ref{sf} and \ref{sg} derive that 
        \begin{align}\label{uu}
            d_{M/\mathcal{F}}\sqrt{\lambda_1^B}\geq \pi+\frac{3}{4}a^2\int_0^{\frac{\pi}{2}}\psi^2(\theta)\d\theta,
        \end{align}where $a=\frac{1-k}{1+k}$.
        
        Therefore,
        \begin{align*}
            \lambda_1^B\geq\frac{\pi^2}{d_{M/\mathcal{F}}^2}.
        \end{align*}

       

        
\end{proof}

\end{document}
